\section*{2. 消灭圈}

\subsection*{2.1. 函子 $R\Psi$}

\subsubsection*{2.1.1.}
设 $S$ 为 0.2.5 中的 Hensel 离散赋值环谱, $\Lambda$ 为挠环, 其挠性与 $S$ 的剩余特征 $p$ 互素（0.2.7）. 更一般地, 可取 $\Lambda$ 为挠性与剩余特征互素的挠层; 最重要的情形是 $\Lambda=\mathbb Z/\ell^n$, 其中 $\ell$ 是在 $\mathcal O_S$ 中可逆的素数.

设 $X$ 为 $S$ 上的概形. 函子 $\Psi$（分别 $\Psi_\eta$）将 $X$（分别 $X_\eta$）上的 $\Lambda$-模层映为 $X_s\times_sS$（分别 $X_s\times_s\eta$）上的 $\Lambda$-模层. 设 $R\Psi$ 和 $R\Psi_\eta$ 为导出函子
\begin{equation*}
R\Psi:D^+(X,\Lambda)\longrightarrow D^+(X_s\times_sS,\Lambda),
\end{equation*}
\begin{equation*}
R\Psi_\eta:D^+(X_\eta,\Lambda)
\longrightarrow D^+(X_s\times_s\eta,\Lambda),
\end{equation*}
并设 $R\Phi$ 为复合函子 $\Phi\circ R\Psi$:
\begin{equation*}
R\Phi:D^+(X,\Lambda)
\longrightarrow D^+(X_s\times_s\eta,\Lambda).
\end{equation*}

\subsubsection*{2.1.2.}
考虑交换图 (1.3.2.1)
\begin{equation}\tag{1.3.2.1}
\begin{tikzcd}
X_{\bar s} \arrow[r, hook, "\bar i"] \arrow[d]
  & \bar X \arrow[d]
  & X_{\bar\eta} \arrow[l, hook', "\bar j"'] \arrow[d] \\
X_s \arrow[r, hook, "i"] & X & X_\eta \arrow[l, hook', "j"']
\end{tikzcd}
\end{equation}
若 $F$ 为 $X$ 上的内射模层, 则其在开子概形 $X_\eta$ 上的限制仍内射. 若 $F_\eta$ 在 $X_\eta$ 上内射, 则其逆像 $F_{\bar\eta}$ 无上同调. 因此 (1.3.3.1) 和 (1.3.3.2) 的导出形式为
\begin{equation}\tag{2.1.2.1}
R\Psi(K)_s=i^*K
\qquad(K\in\operatorname{Ob}D^+(X,\Lambda)),
\end{equation}
\begin{equation}\tag{2.1.2.2}
R\Psi(K)_\eta=R\Psi_\eta(K_\eta)
\qquad(K\in\operatorname{Ob}D^+(X,\Lambda)),
\end{equation}
及
\begin{equation}\tag{2.1.2.3}
R\Psi_\eta(K)_{\bar\eta}
=\bar i^*R\bar j_*K_{\bar\eta}
\qquad(K\in\operatorname{Ob}D^+(X_\eta,\Lambda)).
\end{equation}
这使我们能够不致冲突地记
\begin{equation*}
R\Psi_\eta(K)=(R\Psi(K))_\eta,
\qquad
R\Psi_{\bar\eta}(K)=(R\Psi(K))_{\bar\eta},
\qquad\text{及}\qquad
R\Psi_{\bar\eta}(K)=(R\Psi_\eta(K))_{\bar\eta}.
\end{equation*}
在这些记号下, 区别三角 (1.4.2.2) 成为
\begin{equation}\tag{2.1.2.4}
\begin{tikzcd}[row sep=large, column sep=huge]
{} & R\Phi(K) \arrow[dl] & {} \\
\mathrm{sp}^*i^*K \arrow[rr] & {}
  & R\Psi_\eta(K_\eta) \arrow[ul]
\end{tikzcd}
.\end{equation}
原印本右下项为 $K_s$; 由定义等式 (2.1.2.2), 这里须为 $K_\eta$, 如上所写.

对 $\sigma\in I$, 有\underline{变差}态射
\begin{equation}\tag{2.1.2.5}
\operatorname{Var}(\sigma):R\Phi(K)_{\bar\eta}
\longrightarrow R\Psi_{\bar\eta}(K).
\end{equation}
层 $R^i\Phi(\Lambda)$, 或更一般的 $R^i\Phi(K)$, 称为\underline{消灭圈层}.

\subsubsection*{2.1.3.}
设 $(\bar x,\bar\eta)$ 为由 $X_{\bar s}$ 的一个几何点 $\bar x$ 定义的 $X_s\times_s\eta$ 的几何点（1.2.5）. 我们用 (2.1.2.3) 计算 $R\Psi_\eta(K_\eta)$ 在 $(\bar x,\bar\eta)$ 处的纤维（参见 I, 2.3）. $X$ 在 $\bar x$ 处的严格 Hensel 化 $X_{(\bar x)}$ 是 $S$ 的严格 Hensel 化 $S_{\mathrm{nr}}$ 上的概形, 后者的泛几何点为 $\eta_{\mathrm{nr}}$.

\subsubsection*{命题 2.1.4.}
\underline{有}
\begin{equation*}
(R\Psi(K_\eta))_{(\bar x,\bar\eta)}
=R\Gamma\!\left(
X_{(\bar x)}\times_{\eta_{\mathrm{nr}}}\bar\eta,K
\right);
\end{equation*}
\underline{特别地,}
\begin{equation*}
R^i\Psi(K_\eta)_{(\bar x,\bar\eta)}
=\mathbb H^i\!\left(
X_{(\bar x)}\times_{\eta_{\mathrm{nr}}}\bar\eta,K
\right).
\end{equation*}

应用 SGA~4, XV, 2.1（光滑态射的局部无上同调性定理）, 得到下述重述.

\subsubsection*{重述 2.1.5.}
\underline{若 $f$ 光滑且 $F$ 为局部常值层, 则}
\begin{equation*}
R\Phi(F)=0.
\end{equation*}
更一般地, 若 $K\in\operatorname{Ob}D^+(X,\Lambda)$, 则在 $f$ 光滑且各层 $\mathcal H^i(K)$ 局部常值之处, 有 $R\Phi(K)=0$. 另见 2.1.11.

\subsubsection*{2.1.6.}
设 $f:Y\longrightarrow Y'$ 为 $s$ 上的概形态射. 函子 $f^*$, $f_*$, $f^!$ 和 $f_!$（1.2.7--1.2.9）在导出范畴中有自然的对应物.
\begin{enumerate}[label=\alph*)]
\item 函子 $f^*$ 正合, 因而其导出是平凡的.

\item 函子 $f_*$ 有右导出函子 $Rf_*$. 当 $f$ 拟紧时, 有
\begin{equation}\tag{2.1.6.1}
(Rf_*K)_{\bar s}=Rf_{\bar s*}(K_{\bar s})
\end{equation}
及
\begin{equation}\tag{2.1.6.2}
(Rf_*K)_{\bar\eta}=Rf_{\bar s*}(K_{\bar\eta}).
\end{equation}

\item 当 $f$ 拟有限时, 函子 $f_!$ 正合, 因而其导出是平凡的. 一般说来, $f_!$ 的导出函子具有病态, 须按 SGA~4, XVII 的方式处理. 假设 $f$ 分离且有限型, 并且 $Y'$ 为 Noether 概形. 则存在 [2] 分解 $f=\widetilde f\,j$, 其中 $j$ 为开浸入, $\widetilde f$ 固有, 并定义
\begin{equation*}
Rf_!=R\widetilde f_*\circ j_!.
\end{equation*}
SGA~4, XVII 中的论证说明, $Rf_!$ 不依赖于分解 $\widetilde f\,j$ 的选择. (2.1.6.1) 和 (2.1.6.2) 的类似公式成立.

\item 我们只对拟有限态射使用 $Rf^!$, 此时它就是 $f^!$ 的导出函子. (2.1.6.1) 和 (2.1.6.2) 的类似公式仍然成立.
\end{enumerate}
